\documentclass[CJK]{beamer}
\usepackage{CJKutf8}
\usepackage{beamerthemesplit}
\usetheme{Malmoe}
\useoutertheme[footline=authortitle]{miniframes}
\usepackage{amsmath}
\usepackage{amssymb}
\usepackage{graphicx}
\usepackage{eufrak}
\usepackage{color}
\usepackage{slashed}
\usepackage{simplewick}
\usepackage{tikz}
\usepackage{tcolorbox}
\usepackage{ulem}
\graphicspath{{../figures/}}
%%figures
\def\lfig#1#2{\includegraphics[width=#1 in]{#2}}
\def\tfig#1#2{\includegraphics[height=#1 in]{#2}}
\def\addfig#1#2{\begin{center}\includegraphics[width=#1 in]{#2}\end{center}}
\def\wulian{\includegraphics[width=0.18in]{emoji_wulian.jpg}}
\def\bigwulian{\includegraphics[width=0.35in]{emoji_wulian.jpg}}
\def\bye{\includegraphics[width=0.18in]{emoji_bye.jpg}}
\def\bigbye{\includegraphics[width=0.35in]{emoji_bye.jpg}}
\def\huaixiao{\includegraphics[width=0.18in]{emoji_huaixiao.jpg}}
\def\bighuaixiao{\includegraphics[width=0.35in]{emoji_huaixiao.jpg}}
\def\jianxiao{\includegraphics[width=0.18in]{emoji_jianxiao.jpg}}
\def\bigjianxiao{\includegraphics[width=0.35in]{emoji_jianxiao.jpg}}
\def\haoqi{\includegraphics[width=0.18in]{emoji_haoqi.jpg}}
%% colors
\def\blacktext#1{{\color{black}#1}}
\def\bluetext#1{{\color{blue}#1}}
\def\redtext#1{{\color{red}#1}}
\def\darkbluetext#1{{\color[rgb]{0,0.2,0.6}#1}}
\def\skybluetext#1{{\color[rgb]{0.2,0.7,1.}#1}}
\def\cyantext#1{{\color[rgb]{0.,0.5,0.5}#1}}
\def\greentext#1{{\color[rgb]{0,0.7,0.1}#1}}
\def\darkgray{\color[rgb]{0.2,0.2,0.2}}
\def\lightgray{\color[rgb]{0.6,0.6,0.6}}
\def\gray{\color[rgb]{0.4,0.4,0.4}}
\def\blue{\color{blue}}
\def\red{\color{red}}
\def\orange{\color[rgb]{1.,0.8,0.}}
\def\green{\color{green}}
\def\darkgreen{\color[rgb]{0,0.4,0.1}}
\def\darkblue{\color[rgb]{0,0.2,0.6}}
\def\skyblue{\color[rgb]{0.2,0.7,1.}}
%%control
\def\bral{\left(\begin{array}{l}}
\def\brar{\end{array}\right)}
\def\brall{\left(\begin{array}{ll}}
\def\brarr{\end{array}\right)}
\def\bralll{\left(\begin{array}{lll}}
\def\brarrr{\end{array}\right)}
\def\branchl{\left\{\begin{array}{l}}
\def\branchr{\end{array}\right.}
\def\branchll{\left\{\begin{array}{ll}}
\def\branchrr{\end{array}\right.}
\def\branchlll{\left\{\begin{array}{lll}}
\def\branchrrr{\end{array}\right.}
\def\sfgamma#1{\,\Gamma\left( #1 \right)\,}
\def\be{\begin{equation}}
\def\ee{\nonumber\end{equation}}
\def\bea{\begin{eqnarray}}
\def\eea{\nonumber\end{eqnarray}}
\def\bch{\begin{CJK}{UTF8}{gbsn}}
\def\ech{\end{CJK}}
\def\bitem{\begin{itemize}}
\def\eitem{\end{itemize}}
\def\bcenter{\begin{center}}
\def\ecenter{\end{center}}
\def\bex{\begin{minipage}{0.2\textwidth}\includegraphics[width=0.6in]{jugelizi.png}\end{minipage}\begin{minipage}{0.76\textwidth}}
\def\eex{\end{minipage}}
\def\chtitle#1{\frametitle{\bch#1\ech}}
\def\bmat#1{\left(\begin{array}{#1}}
\def\emat{\end{array}\right)}
\def\bcase#1{\left\{\begin{array}{#1}}
\def\ecase{\end{array}\right.}
\def\bmini#1{\begin{minipage}{#1\textwidth}}
\def\emini{\end{minipage}}
\def\tbox#1{\begin{tcolorbox}#1\end{tcolorbox}}
\def\pfrac#1#2#3{\left(\frac{\partial #1}{\partial #2}\right)_{#3}}
\def\res#1#2{\,\mathrm{res}\,#1\left(#2\right)\,}
\def\newt#1#2{\left(\begin{array}{c}#1\\ #2\end{array}\right)}
\def\reof#1{\,\mathrm{Re}{\left(#1\right)\,}}
\def\imof#1{\,\mathrm{Im}{\left(#1\right)\,}}
\def\Arg#1{\,\mathrm{Arg}\,#1\,}
%%symbols
\def\vleak{\upsilon_{\rm leak}}
\def\cvmol{C_V^{\rm mol}}
\def\cpmol{C_p^{\rm mol}}
\def\intfull{\int_{-\infty}^\infty}
\def\inthalf{\int_0^\infty}
\def\bropt{\,(\ \ \ )}
\def\sone{$\star$}
\def\stwo{$\star\star$}
\def\sthree{$\star\star\star$}
\def\sfour{$\star\star\star\star$}
\def\sfive{$\star\star\star\star\star$}
\def\rint{{\int_\leftrightarrow}}
\def\roint{{\oint_\leftrightarrow}}
\def\stdHf{{\textit{\r H}_f}}
\def\deltaH{{\Delta \textit{\r H}}}
\def\ii{{\dot{\imath}}}
\def\skipline{{\vskip0.1in}}
\def\skiplines{{\vskip0.2in}}
\def\lagr{{\mathcal{L}}}
\def\hamil{{\mathcal{H}}}
\def\vecv{{\mathbf{v}}}
\def\vecx{{\mathbf{x}}}
\def\vecy{{\mathbf{y}}}
\def\veck{{\mathbf{k}}}
\def\vecp{{\mathbf{p}}}
\def\vecn{{\mathbf{n}}}
\def\vecA{{\mathbf{A}}}
\def\vecP{{\mathbf{P}}}
\def\vecsigma{{\mathbf{\sigma}}}
\def\hatJn{{\hat{J_\vecn}}}
\def\hatJx{{\hat{J_x}}}
\def\hatJy{{\hat{J_y}}}
\def\hatJz{{\hat{J_z}}}
\def\hatj#1{\hat{J_{#1}}}
\def\hatphi{{\hat{\phi}}}
\def\hatq{{\hat{q}}}
\def\hatpi{{\hat{\pi}}}
\def\vel{\upsilon}
\def\Dint{{\mathcal{D}}}
\def\adag{{\hat{a}^\dagger}}
\def\bdag{{\hat{b}^\dagger}}
\def\cdag{{\hat{c}^\dagger}}
\def\ddag{{\hat{d}^\dagger}}
\def\hata{{\hat{a}}}
\def\hatb{{\hat{b}}}
\def\hatc{{\hat{c}}}
\def\hatd{{\hat{d}}}
\def\hatD{{\,\hat{D}}}
\def\hatN{{\hat{N}}}
\def\hatH{{\hat{H}}}
\def\hatp{{\hat{p}}}
\def\Fup{{F^{\mu\nu}}}
\def\Fdown{{F_{\mu\nu}}}
\def\newl{\nonumber \\}
\def\vece{\mathrm{e}}
\def\calM{{\mathcal{M}}}
\def\calT{{\mathcal{T}}}
\def\calR{{\mathcal{R}}}
\def\barpsi{\bar{\psi}}
\def\baru{\bar{u}}
\def\barv{\bar{\upsilon}}
\def\qeq{\stackrel{?}{=}}
\def\ftf{\stackrel{\mathcal{FT}}{\Longrightarrow}}
\def\ftb{\stackrel{\mathcal{FT}}{\Longleftarrow}}
\def\ftfb{\stackrel{\mathcal{FT}}{\Longleftrightarrow}}
\def\ltf{\stackrel{\mathcal{LT}}{\Longrightarrow}}
\def\ltb{\stackrel{\mathcal{LT}}{\Longleftarrow}}
\def\ltfb{\stackrel{\mathcal{LT}}{\Longleftrightarrow}}
\def\torder#1{\mathcal{T}\left(#1\right)}
\def\rorder#1{\mathcal{R}\left(#1\right)}
\def\contr#1#2{\contraction{}{#1}{}{#2}#1#2}
\def\trof#1{\mathrm{Tr}\left(#1\right)}
\def\trace{\mathrm{Tr}}
\def\comm#1{\ \ \ \left(\mathrm{used}\ #1\right)}
\def\tcomm#1{\ \ \ (\text{#1})}
\def\slp{\slashed{p}}
\def\slk{\slashed{k}}
\def\calp{{\mathfrak{p}}}
\def\veccalp{\mathbf{\mathfrak{p}}}
\def\Tthree{T_{\tiny \textcircled{3}}}
\def\pthree{p_{\tiny \textcircled{3}}}
\def\dbar{{\,\mathchar'26\mkern-12mu d}}
\def\erf{\mathrm{erf}}
\def\const{\mathrm{const.}}
\def\pheat{\pfrac p{\ln T}V}
\def\vheat{\pfrac V{\ln T}p}
%%units
\def\fdeg{{^\circ \mathrm{F}}}
\def\cdeg{^\circ \mathrm{C}}
\def\atm{\,\mathrm{atm}}
\def\angstrom{\,\text{\AA}}
\def\SIL{\,\mathrm{L}}
\def\SIkm{\,\mathrm{km}}
\def\SIyr{\,\mathrm{yr}}
\def\SIGyr{\,\mathrm{Gyr}}
\def\SIV{\,\mathrm{V}}
\def\SImV{\,\mathrm{mV}}
\def\SIeV{\,\mathrm{eV}}
\def\SIkeV{\,\mathrm{keV}}
\def\SIMeV{\,\mathrm{MeV}}
\def\SIGeV{\,\mathrm{GeV}}
\def\SIcal{\,\mathrm{cal}}
\def\SIkcal{\,\mathrm{kcal}}
\def\SImol{\,\mathrm{mol}}
\def\SIN{\,\mathrm{N}}
\def\SIHz{\,\mathrm{Hz}}
\def\SIm{\,\mathrm{m}}
\def\SIcm{\,\mathrm{cm}}
\def\SIfm{\,\mathrm{fm}}
\def\SImm{\,\mathrm{mm}}
\def\SInm{\,\mathrm{nm}}
\def\SImum{\,\mathrm{\mu m}}
\def\SIJ{\,\mathrm{J}}
\def\SIW{\,\mathrm{W}}
\def\SIkJ{\,\mathrm{kJ}}
\def\SIs{\,\mathrm{s}}
\def\SIkg{\,\mathrm{kg}}
\def\SIg{\,\mathrm{g}}
\def\SIK{\,\mathrm{K}}
\def\SImmHg{\,\mathrm{mmHg}}
\def\SIPa{\,\mathrm{Pa}}
%page
\def\secpage#1#2{\begin{frame}\bch\bcenter{\bf \Huge #1} \skipline \tbox{#2}\ecenter\ech\end{frame}}

\def\courseurl{http://zhiqihuang.top}

\def\tpage#1#2{
\begin{frame}
\begin{center}
\begin{Large}
\bch
热学 \\
第#1讲 #2

{\vskip 0.3in}

黄志琦

\ech
\end{Large}
\end{center}

\vskip 0.2in

\bch
教材：《热学》第二版，赵凯华，罗蔚茵，高等教育出版社
\ech

\bch
课件下载
\ech
\courseurl
\end{frame}
}

\def\bfr#1{
\begin{frame}
\chtitle{#1} 
\bch
}

\def\efr{
\ech 
\end{frame}
}

\title{Lesson 04 - Maxwell Distribution}
  \author{}
  \date{}
\begin{document}
\tpage{4}{麦克斯韦分布}

\section{Reivew}

\begin{frame}
\chtitle{上一讲内容回顾}
\bch
上一讲的内容基本上都不考
\ech
\end{frame}


\begin{frame}
\chtitle{本讲内容预告}
\bch
\bitem
\item{概率密度函数}
\item{高斯积分的技巧}
\item{微观粒子的态}
\item{麦克斯韦分布}
\eitem
\ech
\end{frame}

\section{Probability Density Function}

\begin{frame}
\chtitle{概率密度函数(Probability Density Function)}
\bch
\addfig{3}{pdf.jpg}

请一位同学来解释一下概率密度函数的概念
\ech
\end{frame}

\begin{frame}
\chtitle{概率密度函数的换元}
\bch
假设随机变量$x$的概率密度函数为$P(x)$，我们把$x$以一定的规则映射到某个新的随机变量$y$，并考察$y$的概率密度函数$Q(y)$。

\skiplines

做变量替换后要注意不变量是概率元:
$$Q(y) dy = \sum_{x\rightarrow y} P(x) dx$$
也就是说，把所有能映射为$y$的$x$附近的概率加起来，然后除以$dy$，得到的就是$y$的概率密度函数。
\ech
\end{frame}


\begin{frame}
\chtitle{思考题}
\bch

\addfig{0.5}{think.jpg}

假设随机变量$x$的概率密度函数为$P(x)$，定义$y=x/2$，求$y$的概率密度函数。
\ech
\end{frame}

\begin{frame}
\chtitle{思考题}
\bch

\addfig{0.7}{think2.jpg}

假设随机变量$x$的概率密度函数为$P(x)$，定义$y=x^2$，求$y$的概率密度函数。
\ech
\end{frame}


\begin{frame}
\chtitle{很多个一样的东西的和总是满足高斯分布}
\bch
    {\blue 中心极限定理：满足任何给定分布的多个随机变量的和总是趋向于服从高斯分布（也叫正态分布）:}

    (我认为统计学里只有两条非平凡的定律，中心极限定理是其中一条)

    \skiplines
    
   高斯分布的{\bf 概率密度函数}为：
   $$ P(x) = \frac{1}{\sqrt{2\pi}\sigma} e^{-\frac{(x-\mu)^2}{2\sigma^2}}  $$

   其中$\sigma$是标准差，$\mu$是均值。


\ech
\end{frame}

\begin{frame}
\chtitle{很多个一样的东西的和总是满足高斯分布}
\bch
\bex
张同学测量体重，得到的数值是一个随机变量。这个随机变量服从的分布可能非常复杂（取决于测量工具，测量方式等）。但多次测量的和（或者平均）总是趋向于服从正态分布。
多次测量求平均后的结果可以写成 $(175.0\pm0.5)$ （斤？公斤？吨？）的形式: 默认就是$\mu = 175$, $\sigma = 0.5$的正态分布。
\eex
\ech
\end{frame}


\begin{frame}
\chtitle{把问题简单化}
\bch
因为$\sigma$反映了$x$的大致变化尺度，$\mu$代表了中心值。我们总是可以重新定义变量（相当于重新取零点和单位）
$$y \equiv \frac{x-\mu}{\sigma}$$
\ech
\end{frame}


\section{Microstate}

\begin{frame}
\chtitle{态的例子}
\bch
\bitem
\item{硬币有两种态：正面朝上，反面朝上

\addfig{0.7}{coins.jpg}
}
\item{骰子有六种态：1，2，3，4，5，6

\addfig{0.3}{touzi.png}
}
\item{同学们有很多种态：上课，吃饭，睡觉，和床引力战斗……

\addfig{1.5}{xbxstates.jpg}
}

\eitem
\ech
\end{frame}

\begin{frame}
\chtitle{有的人反对了: 骰子的六种态不一定是123456
}
\bch

\addfig{1.2}{touzi_work.jpg}

\skipline

\wulian好吧，这不是我们今天要讲的重点…
\ech
\end{frame}

\begin{frame}
\chtitle{相空间}
\bch
气体分子的“态”由位置$(x, y, z)$和速度$(\upsilon_x, \upsilon_y, \upsilon_z)$描述：
$$(x, y, z, \upsilon_x, \upsilon_y, \upsilon_z)$$
或者更专业的说法是由位置和动量描述：
$$(x, y, z, \calp_x, \calp_y, \calp_z)$$

由$x, y, z, \calp_x, \calp_y, \calp_z$这六个变量（坐标轴）张成的空间称为{\blue \bf 相空间}。

\bitem
\item{经典图像（错误）：分子的态由相空间中的一个点来描述。}
\item{量子图像（正确）：{\bf \blue 分子的态由相空间中的一个“小块”来描述。}}
\eitem

\ech
\end{frame}


\begin{frame}
\chtitle{相空间里的“小块”是什么鬼？}
\bch
假设位置空间$x$是一维的，则相空间是$(x, \calp)$构成的二维空间，相空间里的“小块”如下图所示:
\bmini{0.35}
\addfig{0.8}{phasespaceblock.png}
\emini
\bmini{0.6}
\bitem
\item{普朗克常数$h=6.63\times 10^{-34}\SIJ\cdot\SIs$}
\item{只要面积为$h$，即是合法的“小块”。如何取$\Delta x$或$\Delta \calp$视测量需要而定。}
\item{量子力学的测不准原理保证了“小块”内的点是无法通过测量区分的。}
\eitem
\emini

\skipline

在三维空间的情况，相空间为六维空间，“小块”体积为$h^3$ (六条边长满足$\Delta \calp_x\Delta x =\Delta \calp_y\Delta y =\Delta \calp_z\Delta z = h$)
\ech
\end{frame}


\begin{frame}
\chtitle{态是离散的相空间小块}
\bch
\bitem
\item{\bf $n$维位置空间$(x_1, x_2, \ldots, x_n)$对应$2n$维相空间$(x_1, x_2, \ldots, x_n, \calp_1, \calp_2, \ldots, \calp_n)$。}
\item{\bf 相空间内“小块”是一个假想的$2n$维体。在每个轴方向上的边长依次为$(\Delta x_1, \Delta x_2, \ldots, \Delta x_n, \Delta \calp_1, \Delta \calp_2, \ldots, \Delta \calp_n)$。边长两两成对满足$\Delta x_i \cdot \Delta \calp_i = h$ ($i=1, 2, \ldots, n$)。显然，“小块”的体积为$h^n$。}
\item{\bf 在每个维度上如何取$\Delta x_i$则视测量需求而定，有时甚至可以取$\Delta x_i$为宏观尺度。}
\item{\bf 相空间的每个“小块”对应一个微观粒子的态。虽然数学上“小块”仍然是无限可分的，在物理上却由于测不准原理而无法再划分更细的状态。}
\eitem
\ech
\end{frame}

\begin{frame}
\chtitle{思考题}
\bch
一个边长为$0.663\SIm$的正方体容器里，质量为$10^{-26}\SIkg$，速度不超过$10^3\SIm\SIs^{-1}$的粒子共有多少种可能的态？
\ech
\end{frame}

\begin{frame}
\chtitle{思考题}
\bch
一个边长为$0.663\SIm$的正方体容器里，质量为$10^{-26}\SIkg$，速度不超过$0.8$倍光速的粒子共有多少种可能的态？
\ech
\end{frame}


\begin{frame}
\chtitle{测量空气分子的动量}
\bch
考虑装在一个边长$L= 1\SIm$的正方体容器里的空气。取空气分子的平均质量，$\bar{m} = \frac{29\SIg\cdot\SImol^{-1}}{6.02\times 10^{23}\SImol^{-1}}\approx 5 \times 10^{-26}\SIkg$

取室温$T=300K$，则空气分子动量大小的数量级为
$$\calp \sim \sqrt{m^2\overline{\upsilon^2}} \sim \sqrt{3mkT} \approx 2\times 10^{-23} \SIkg \SIm\SIs^{-1} $$

我们准备测量空气分子的速度的分布规律，希望动量空间的分辨率尽可能地高，所以在位置空间尽可能地降低分辨率：取$\Delta x=\Delta y=\Delta z = L$（也就是说我们只要求分子在容器内，而不去测量它的具体位置）。那么动量空间的分辨率(即小块沿动量轴方向的边长)为
$$\Delta \calp_x = \Delta \calp_y = \Delta \calp_z= \frac{h}{L} = 6.63\times 10^{-34} \SIkg\SIm\SIs^{-1} \sim 10^{-11} \calp$$

\ech
\end{frame}


\begin{frame}
\chtitle{测量空气分子的动量(续）}
\bch
现在我们更贪心，希望把分子的位置确定到分子平均距离的数量级（$d\sim 3\times 10^{-9}\SIm$)，这时动量分辨率约为
$$\Delta \calp \sim \frac{h}{d} \sim 10^{-25} \SIkg\SIm\SIs^{-1} \sim 10^{-2} \calp$$

也就是说，课本上第二章讨论玻尔兹曼分布时的经典图像（同时知道分子的位置和动量），大致上可以实现。

\ech
\end{frame}


\section{Maxwell Distribution}

\begin{frame}
\chtitle{麦克斯韦分布(Maxwell distribution)}
\bch
现在我们考虑非相对论气体分子的速度分布规律。

\skipline

{\blue \bf 分子的“态” = 分子在相空间的哪个小块里 }

\skipline

我们仍取$\Delta x$, $\Delta y$, $\Delta z$为整个容器的长，宽，高。那么在位置空间就不需要进行划分了。分子的态由动量空间的坐标$(\calp_x, \calp_y, \calp_z)$唯一决定。当然，$\calp_x$, $\calp_y$, $\calp_z$的取值并不是连续的，它们分别是$\Delta \calp_x=\frac{h}{\Delta x}$, $\Delta \calp_y=\frac{h}{\Delta y}$, $\Delta \calp_z=\frac{h}{\Delta z}$的整数倍。

\skipline
根据我们第一讲给的“万能法则”，达到热平衡时，分子处在一个能量为$\varepsilon$的态的概率正比于$e^{-\frac{\varepsilon}{kT}}$。
所以，质量为$m$的分子处在任何一个态$(\calp_x, \calp_y, \calp_z)$的概率正比于
$$e^{-\frac{m\upsilon^2}{2kT}}$$
其中
$\upsilon = \frac{\sqrt{\calp_x^2 +\calp_y^2+ \calp_z^2}}{m}$
是分子速率。
\ech
\end{frame}

\begin{frame}
\chtitle{麦克斯韦分布(续)}
\bch
分子速率在$\upsilon$和$\upsilon+d\upsilon$之间的态（相空间“小块”）有多少个呢？很简单，用相空间的体积除以小块的体积$h^3$即可。

\skipline


态的数目 = $ \frac{(\Delta x \Delta y \Delta z )4\pi (m\upsilon)^2d(m\upsilon)}{h^3}\propto \upsilon^2d\upsilon$
\ech
\end{frame}

\begin{frame}
\chtitle{麦克斯韦分布(续)}
\bch
{\small 分子速率在$\upsilon$和$\upsilon+d\upsilon$之间的概率=每个态上出现的概率$\times$态的数目}，
$$F_M(\upsilon) d\upsilon \propto \upsilon^2 e^{-\frac{m\upsilon^2}{2kT}}d\upsilon.$$
上式中的$F_M(\upsilon)$是分子速率为$\upsilon$的概率密度，满足归一化条件$\int_0^\infty F_M(\upsilon)d\upsilon = 1$。在三维速度空间，分子速度为$\vecv$(注意这是个三维矢量)的概率密度
$$f_M(\vecv) = \frac{1}{4\pi|\vecv|^2}F_M(|\vecv|)\, ,$$
满足归一化条件$\int f_M(\vecv) d^3\vecv = 1$ (积分元$d^3\vecv \equiv d\upsilon_xd\upsilon_yd\upsilon_z$).

教材上有时把$f_M(\vecv)$写成$f_M(\upsilon)$, $f_M(\upsilon^2)$,$f(\upsilon)$等等，你们当作没看见就行\wulian

\ech
\end{frame}

\begin{frame}
\chtitle{麦克斯韦分布的归一化系数}
\bch
设
$$F_M(\upsilon) = C \upsilon^2 e^{-\frac{m\upsilon^2}{2kT}}.$$
由归一化条件$\int_0^\infty F_M(\upsilon)d\upsilon = 1$得到
$$ C = \frac{1}{\int_0^\infty \upsilon^2 e^{-\frac{m\upsilon^2}{2kT}}d\upsilon} \, . $$
下面我们来计算这个积分。
\ech
\end{frame}


\begin{frame}
\chtitle{归一化的麦克斯韦速度分布函数}
\bch
{\small
于是我们得到
$$F_M(\upsilon) = 4\pi \upsilon^2 \left(\frac{m}{2\pi kT}\right)^{3/2} e^{-\frac{m\upsilon^2}{2kT}}\,.$$

写成{\bf \blue 麦克斯韦速度分布函数（即三维速度空间的概率密度）：
$$f_M(\vecv) = \left(\frac{m}{2\pi kT}\right)^{3/2} e^{-\frac{m|\vecv|^2}{2kT}}\, .$$}
}
\ech
\end{frame}

\begin{frame}
\chtitle{另一种推导方法(更方便快捷)}
\bch
{\small
由于动能在三个自由度上的贡献是独立的，我们由“万能法则”可以得到$\upsilon_x$的概率密度函数为
$$ f_{1D}(\upsilon_x) =  C e^{-\frac{m\upsilon_x^2}{2kT}}$$
由归一化条件$\int_{-\infty}^\infty f(\upsilon_x)d\upsilon_x = 1$以及高斯积分公式($a = \frac{m}{2kT}$)，得到

$$ f_{1D}(\upsilon_x) =   \sqrt{\frac{m}{2\pi kT}} e^{-\frac{m\upsilon_x^2}{2kT}}$$
显然对$\upsilon_y$, $\upsilon_z$可以得到一样的结果，因为综合事件的概率等于每一个独立子事件的概率的乘积，所以对$\vecv = (\upsilon_x, \upsilon_y, \upsilon_z)$有
$$f_M(\vecv) = f_{1D}(\upsilon_x)f_{1D}(\upsilon_y)f_{1D}(\upsilon_z)=  \left(\frac{m}{2\pi kT}\right)^{3/2} e^{-\frac{m(\upsilon_x^2+\upsilon_y^2+\upsilon_z^2)}{2kT}}   \, .$$
}
\ech
\end{frame}

\begin{frame}
\chtitle{统计平均量}
\bch
速度的任何函数$g(\vecv)$的统计平均显然为
{\blue $$\overline{g(\vecv)} = \int f_M(\vecv) g(\vecv) d^3\vecv$$}
如果$g$只是速率的函数，则其统计平均可以写成
{\blue $$\overline{g(\upsilon)} = \int_0^{\infty} F_M(\upsilon) g(\upsilon)d\upsilon.$$}
如果$g$只是单独一个自由度$\upsilon_x$的函数，则其统计平均
{\blue $$\overline{g(\upsilon_x)} = \int_{-\infty}^{\infty} f_{1D}(\upsilon_x) g(\upsilon_x)d\upsilon_x.$$}

\ech
\end{frame}

\begin{frame}
\chtitle{练习题}
\bch
{\small
按麦克斯韦分布的三维概率密度形式：
$$f_M(\upsilon_x, \upsilon_y, \upsilon_z) =\left(\frac{m}{2\pi kT}\right)^{3/2} e^{-\frac{m(\upsilon_x^2+\upsilon_y^2+\upsilon_z^2)}{2kT}}   \, .$$
或者其径向一维概率密度形式：
$$F_M(\upsilon) = 4\pi \upsilon^2 \left(\frac{m}{2\pi kT}\right)^{3/2} e^{-\frac{m\upsilon^2}{2kT}}\,\ \ \ \ (\upsilon>0)$$
或其单独一个自由度的一维概率密度形式：
$$ f_{1D}(\upsilon_x) =   \sqrt{\frac{m}{2\pi kT}} e^{-\frac{m\upsilon_x^2}{2kT}}$$

计算下列函数的统计平均量
\bitem
\item{ $\overline{\upsilon}$ }
\item{ $\overline{\upsilon^2}$ }
\item{ $\overline{\upsilon_x\theta(\upsilon_x)}$，其中台阶函数$\theta(x)$当$x\ge 0$时取值为$1$,否则为零。}
\eitem}
\ech
\end{frame}


\begin{frame}
\chtitle{方均根速率(root mean square velocity)}
\bch
{\small 方均根速率$\upsilon_{\rm rms}$定义为
$$\upsilon_{\rm rms} \equiv \sqrt{\overline{\upsilon^2}}$$
它表征的是分子平均动能的大小。麦克斯韦分布的方均根速率为$\upsilon_{\rm rms} = \sqrt{\frac{3kT}{m}}$，证明如下：}
{\tiny
\bmini{0.6}
\bea
\overline{\upsilon^2} &=& \int F_M(\upsilon) \upsilon^2d\upsilon \newl
&=& 4\pi  \left(\frac{m}{2\pi kT}\right)^{3/2} \int_0^\infty \upsilon^4 e^{-\frac{m\upsilon^2}{2kT}} d\upsilon \newl
&=& 4\pi  \left(\frac{\alpha}{\pi}\right)^{3/2}\left. \frac{d^2}{d\alpha^2}\left(\int_0^\infty e^{-\alpha \upsilon^2}d\upsilon \right)\right\vert_{\alpha = m/(2kT)} \newl
&=& 4\pi  \left(\frac{\alpha}{\pi}\right)^{3/2} \left. \frac{d^2}{d\alpha^2}\left(\frac{1}{2}\sqrt{\frac{\pi}{\alpha}}\right)\right\vert_{\alpha = m/(2kT)} \newl
&=& 4\pi  \left.\left(\frac{\alpha}{\pi}\right)^{3/2} \frac{3}{8\alpha^2} \sqrt{\frac{\pi}{\alpha}} \right\vert_{\alpha = m/(2kT)} \newl
&=& \frac{3kT}{m} 
\eea 
\emini
\bmini{0.35}
把$\upsilon^4$看成是两次求导的产物，并心安理得地交换了积分和求导次序

\skiplines

积分范围是$0$到$\infty$,所以是高斯积分的一半。
\vspace{0.05in}
\emini
}
\ech
\end{frame}


\begin{frame}
\chtitle{平均速率(mean velocity)}
\bch
{\small 
平均速率定义为所有分子的速率的平均值$$\overline{\upsilon}$$
它可以用来计算分子的平均自由程。麦克斯韦分布的平均速率为$\bar{\upsilon} = \sqrt{\frac{8kT}{\pi m}}$，证明如下:}
{\tiny
\bmini{0.6}
\bea
\bar{\upsilon} &=& \int_0^\infty F_M(\upsilon)\upsilon d\upsilon \newl
&=& 4\pi  \left(\frac{m}{2\pi kT}\right)^{3/2} \int_0^\infty \upsilon^3 e^{-\frac{m\upsilon^2}{2kT}} d\upsilon \newl
&=& 4\pi  \left(\frac{\alpha}{\pi}\right)^{3/2}\left. \frac{d}{d\alpha}\left(\int_0^\infty (-\upsilon e^{-\alpha \upsilon^2})d\upsilon \right)\right\vert_{\alpha = m/(2kT)} \newl
&=& 4\pi  \left(\frac{\alpha}{\pi}\right)^{3/2}\left. \frac{d}{d\alpha}\left(\left.\frac{1}{2\alpha} e^{-\alpha \upsilon^2}\right\vert_0^\infty \right)\right\vert_{\alpha = m/(2kT)} \newl
&=& 4\pi  \left(\frac{\alpha}{\pi}\right)^{3/2}\left. \frac{d}{d\alpha}\left(-\frac{1}{2\alpha} \right)\right\vert_{\alpha = m/(2kT)} \newl
&=& \sqrt{\frac{8kT}{\pi m}}
\eea
\emini
\bmini{0.35}
把$(-\upsilon^2)$看成是一次求导的产物，交换了积分和求导次序

\vspace{0.3in}
\emini
}
\ech
\end{frame}


\begin{frame}
\chtitle{泄流速率}
\bch
{\small
泄流速率定义为
$$\overline{\upsilon_x^+} \equiv \overline{\upsilon_x\theta(\upsilon_x)}$$
设分子数密度为$n$，在容器壁上挖个小孔，则单位时间从单位面积上泄出的分子数为
$n \overline{\upsilon_x^+}$。麦克斯韦分布下的泄流速率为$\sqrt{\frac{kT}{2\pi m}}$，证明如下：}
{\tiny
\bea
\overline{\upsilon_x^+} &=& \int_{-\infty}^\infty f_{1D}(\upsilon_x) \upsilon_x \theta(\upsilon_x) d\upsilon_x \newl
 &=& \int_0^\infty f_{1D}(\upsilon_x) \upsilon_x  d\upsilon_x \newl
 &=& \int_0^\infty \sqrt{\frac{m}{2\pi kT}} e^{-\frac{m\upsilon_x^2}{2kT}}\upsilon_x d\upsilon_x \newl
&=& \sqrt{\frac{m}{2\pi kT}} \left(-\frac{kT}{m}\right) \left. e^{-\frac{m\upsilon_x^2}{2kT}}\right\vert_0^{\infty} \newl
&=& \sqrt{\frac{kT}{2\pi m}}
\eea
}

\ech
\end{frame}

\begin{frame}
\chtitle{思考题}
\bch
\addfig{0.8}{think3.jpg}
“泄流”和“泻流”哪个词里的汉字用法更为正确？
\ech
\end{frame}


\section{MB distribution}

\begin{frame}
\chtitle{麦克斯韦-玻尔兹曼分布(Maxwell-Boltzmann Distribution)}
\bch
\bitem
\item{如果分子在不同位置有不同的势能，则分子在位置空间的分布就不是均匀的了。也就是说，气体各处会产生密度梯度。}
\item{如果我们希望同时描述分子在位置空间和速度空间的分布规律，划分相空间“小块”时就要两者兼顾了: 即$\Delta x$和
$\Delta \calp_x$都必须远小于我们测量的精度，才能用连续的概率密度函数来描述分子的分布。}
\eitem
\ech
\end{frame}


\begin{frame}
\chtitle{重力势能的情况}
\bch
假设沿$z$方向有强度为$g$的重力场，则分子的能量
$$\varepsilon = \frac{1}{2}m\upsilon^2 + mgz$$

按照“万能法则”，分子在相空间某小块(设坐标为$(x, y, z, \calp_x, \calp_y, \calp_z)$）出现的概率正比于：

$$f_{MB}(x, y, z, \upsilon_x, \upsilon_y,\upsilon_z) \propto e^{-\frac{\frac{1}{2}m\upsilon^2 + mgz}{kT}}$$

上式可以拆成两个独立分布的乘积：
$$f_{MB}(x, y, z, \upsilon_x, \upsilon_y,\upsilon_z) \propto e^{-\frac{m\upsilon^2}{2kT}} e^{-\frac{mgz}{kT}}$$

可以看到，在速度空间的分布和在位置空间的分布互不干扰，所以我们可以得到分子数密度
$$n = n_0 e^{-\frac{mgz}{kT}}\, ,$$
其中$n_0$为$z=0$处的分子数密度。
\ech
\end{frame}

\begin{frame}
\chtitle{离心力势能}
\bch
在旋转参照系中，离心力势能可以写成
$$U(r) = -\frac{1}{2}m\omega^2r^2$$
其中$r$为离旋转中心的径向距离，$\omega$为旋转系相对于惯性系的旋转角速度。

\skipline

根据一样的推导过程，可以得到
$$n = n_0 e^{\frac{m\omega^2r^2}{2kT}}$$
\ech
\end{frame}


\begin{frame}
\chtitle{思考题}
\bch
\addfig{0.8}{think.jpg}
你能否解释台风的“外围狂风暴雨，中心风和日丽”的奇特现象？
\ech
\end{frame}

\begin{frame}
\chtitle{麦克斯韦-玻尔兹曼能量分布律}
\bch
一般地，如果位置空间和速度空间的能量形式没有交叉项，我们可以分离位置空间和速度空间的分布（概率密度），得到分子在相空间的概率密度为：
$$f_{MB}(x, y, z,\upsilon_x, \upsilon_y,\upsilon_z) = n_0 \left(\frac{m}{2\pi kT}\right)^{3/2} e^{-\frac{\varepsilon}{kT}} $$
其中能量$\varepsilon$是动能与势能之和：
$$\varepsilon = \frac{1}{2}m\upsilon^2 + U(x,y,z)$$
这是著名的麦克斯韦-玻尔兹曼能量分布律（\wulian 好像跟我们反复念叨的万能法则是一回事）。
\ech
\end{frame}


\begin{frame}
\chtitle{下周内容预告}
\bch

下周我们将解释万能法则的由来
\ech
\end{frame}

\begin{frame}
\chtitle{第三周作业(序号接第二周）}
\bch

\bitem
\item[7]{教材习题2-2}
\item[8]{教材习题2-3}
\item[9]{教材习题2-19}
\eitem

\ech
\end{frame}

\end{document}
